% Options for packages loaded elsewhere
\PassOptionsToPackage{unicode}{hyperref}
\PassOptionsToPackage{hyphens}{url}
\documentclass[
  ignorenonframetext,
  aspectratio=169,
]{beamer}
\newif\ifbibliography
\usepackage{pgfpages}
\setbeamertemplate{caption}[numbered]
\setbeamertemplate{caption label separator}{: }
\setbeamercolor{caption name}{fg=normal text.fg}
\beamertemplatenavigationsymbolsempty
% remove section numbering
\setbeamertemplate{part page}{
  \centering
  \begin{beamercolorbox}[sep=16pt,center]{part title}
    \usebeamerfont{part title}\insertpart\par
  \end{beamercolorbox}
}
\setbeamertemplate{section page}{
  \centering
  \begin{beamercolorbox}[sep=12pt,center]{section title}
    \usebeamerfont{section title}\insertsection\par
  \end{beamercolorbox}
}
\setbeamertemplate{subsection page}{
  \centering
  \begin{beamercolorbox}[sep=8pt,center]{subsection title}
    \usebeamerfont{subsection title}\insertsubsection\par
  \end{beamercolorbox}
}
% Prevent slide breaks in the middle of a paragraph
\widowpenalties 1 10000
\raggedbottom
\AtBeginPart{
  \frame{\partpage}
}
\AtBeginSection{
  \ifbibliography
  \else
    \frame{\sectionpage}
  \fi
}
\AtBeginSubsection{
  \frame{\subsectionpage}
}
\usepackage{iftex}
\ifPDFTeX
  \usepackage[T1]{fontenc}
  \usepackage[utf8]{inputenc}
  \usepackage{textcomp} % provide euro and other symbols
\else % if luatex or xetex
  \usepackage{unicode-math} % this also loads fontspec
  \defaultfontfeatures{Scale=MatchLowercase}
  \defaultfontfeatures[\rmfamily]{Ligatures=TeX,Scale=1}
\fi
\usepackage{lmodern}
\ifPDFTeX\else
  % xetex/luatex font selection
\fi
% Use upquote if available, for straight quotes in verbatim environments
\IfFileExists{upquote.sty}{\usepackage{upquote}}{}
\IfFileExists{microtype.sty}{% use microtype if available
  \usepackage[]{microtype}
  \UseMicrotypeSet[protrusion]{basicmath} % disable protrusion for tt fonts
}{}
\makeatletter
\@ifundefined{KOMAClassName}{% if non-KOMA class
  \IfFileExists{parskip.sty}{%
    \usepackage{parskip}
  }{% else
    \setlength{\parindent}{0pt}
    \setlength{\parskip}{6pt plus 2pt minus 1pt}}
}{% if KOMA class
  \KOMAoptions{parskip=half}}
\makeatother
\usepackage{longtable,booktabs,array}
\usepackage{caption}
\captionsetup[table]{skip=6pt}
\newcounter{none} % for unnumbered tables
\usepackage{calc} % for calculating minipage widths
\usepackage{caption}
% Make caption package work with longtable
\makeatletter
\def\fnum@table{\tablename~\thetable}
\makeatother
\setlength{\emergencystretch}{3em} % prevent overfull lines
\providecommand{\tightlist}{%
  \setlength{\itemsep}{0pt}\setlength{\parskip}{0pt}}
% Auto-generated by infrastructure/rendering/_pdf_latex_helpers.py::extract_math_font_preamble.
% Minimal header passed to Pandoc via -H so Beamer slides pick up the same math font
% as the combined-PDF path. Adjust the manuscript's preamble.md to change the font globally.
\usepackage{unicode-math}
\setmathfont{latinmodern-math.otf}

% Manuscript macro/package fallbacks (newcommand -> providecommand).
% Core mathematics
\usepackage{amsmath}
\usepackage{amssymb}
% Document layout
% Tables
\usepackage{booktabs}
% Caption styling (booktabs already loaded above). Small caption font with a
% bold label ("Figure 1", "Table 2") gives the math/figure-heavy layout a
% consistent, journal-like caption rhythm without fighting pandoc-crossref —
% pandoc emits standard \caption{}, and caption/captionsetup only restyle it.
% ── Theorem environments ─────────────────────────────────────────────────
% amsthm is loaded above. The FedGVI<->Friston bridge is stated as a small
% number of theorems/definitions/lemmas/propositions/corollaries; sharing one
% counter (definition/lemma/proposition/corollary all step the `theorem`
% counter) gives a single monotone numbering 1, 2, 3, ... across all kinds, so
% authors can reference them as Theorem 1, Definition 2, Corollary 3 without a
% per-kind counter drifting out of order. Plain style for theorem-like results,
% definition style (upright body) for definitions.
% (algorithm2e intentionally not loaded: this manuscript states results as
% numbered amsthm Theorems/Definitions, not algorithm2e pseudocode.)
% Code listings
\usepackage{listings}
% Typography and formatting
% (siunitx intentionally not loaded: no \SI/\num units appear in this manuscript.)
% Cross-references and citations
% (cleveref and titlesec intentionally not loaded: unavailable in this TeX tree
% and not required. Cross-references resolve through pandoc-crossref's own
% \ref-style names — no \cref appears — and section heading styling falls back
% to the pandoc/LaTeX defaults.)
% ── TeX Live Basic-compatible mono font for code listings ────────────
% Latin Modern Mono is bundled with TeX Live Basic and keeps the manuscript
% renderable on minimal TeX installations. Mathematical Unicode coverage is
% provided separately by Latin Modern Math below, so the code font does not
% need a private JuliaMono installation.
%
% Use the TeX font filename rather than the system-family name: the minimal
% TeX Live fontconfig database may not register the OpenType family even though
% kpsewhich can resolve the bundled files.
% Math font for unicode-math: Latin Modern Math (TeX Live) has full BMP
% coverage including U+2223 (\mid), U+226A/226B (\ll/\gg), and the Greek/
% blackboard letters used in equations. Without an explicit \setmathfont,
% unicode-math falls back to lmroman text font which lacks several glyphs
% and emits "Missing character" warnings on every \mid in math mode.

% Slide layout defaults for warning-clean scientific decks.
\usepackage{etoolbox}
\IfFileExists{xurl.sty}{\usepackage{xurl}}{}
\IfFileExists{seqsplit.sty}{\usepackage{seqsplit}}{\newcommand{\seqsplit}[1]{#1}}
\protected\def\breakseq#1{\seqsplit{#1}}
\protected\def\breaktt#1{\begingroup\ttfamily\seqsplit{#1}\endgroup}
\setlength{\emergencystretch}{6em}
\tolerance=5000
\hbadness=10000
\hfuzz=1pt
\setlength{\tabcolsep}{2pt}
\AtBeginEnvironment{longtable}{\tiny\renewcommand{\arraystretch}{0.86}\setlength{\tabcolsep}{1pt}}
\AtBeginEnvironment{tabular}{\tiny\renewcommand{\arraystretch}{0.86}\setlength{\tabcolsep}{1pt}}
\AtBeginEnvironment{equation}{\tiny}
\AtBeginEnvironment{equation*}{\tiny}
\AtBeginEnvironment{align}{\tiny}
\AtBeginEnvironment{align*}{\tiny}
\AtBeginEnvironment{itemize}{\footnotesize}
\AtBeginEnvironment{enumerate}{\footnotesize}
\AtBeginEnvironment{description}{\footnotesize}
\setbeamerfont{caption}{size=\tiny}
\setbeamerfont{caption name}{size=\tiny}
\setbeamerfont{normal text}{size=\small}
\setbeamerfont{frametitle}{size=\small}
\setbeamerfont{section title}{size=\footnotesize}
\setbeamerfont{subsection title}{size=\footnotesize}
\setbeamertemplate{section page}{%
  \centering
  \begin{beamercolorbox}[sep=12pt,center,wd=\paperwidth]{section title}
    \parbox{0.86\paperwidth}{\centering\usebeamerfont{section title}\insertsection\par}
  \end{beamercolorbox}
}
\setbeamertemplate{subsection page}{%
  \centering
  \begin{beamercolorbox}[sep=8pt,center,wd=\paperwidth]{subsection title}
    \parbox{0.86\paperwidth}{\centering\usebeamerfont{subsection title}\insertsubsection\par}
  \end{beamercolorbox}
}
\setlength{\abovecaptionskip}{2pt}
\setlength{\belowcaptionskip}{0pt}

% Natbib and cross-reference fallbacks — slides don't load natbib
% or cleveref, but combined-PDF manuscript prose may emit these
% commands. The fallback renders citations as a bracketed key list
% and cross-references as detokenized labels so slides don't fail on
% undefined control sequences or raw underscores. \providecommand is
% a no-op if packages load later.
\providecommand{\citep}[1]{[#1]}
\providecommand{\citet}[1]{#1}
\providecommand{\citealp}[1]{#1}
\providecommand{\citeauthor}[1]{#1}
\providecommand{\citeyear}[1]{#1}
\providecommand{\cref}[1]{\texttt{\detokenize{#1}}}
\providecommand{\Cref}[1]{\texttt{\detokenize{#1}}}
\providecommand{\crefrange}[2]{\texttt{\detokenize{#1}}--\texttt{\detokenize{#2}}}
\providecommand{\Crefrange}[2]{\texttt{\detokenize{#1}}--\texttt{\detokenize{#2}}}
\providecommand{\paragraph}[1]{\textbf{#1}\ }

% Opt-in accessible presentation profile.
\setbeamerfont{normal text}{size*={20pt}{24pt}}
\setbeamerfont{frametitle}{size*={28pt}{32pt}}
\setbeamerfont{section title}{size*={28pt}{32pt}}
\setbeamerfont{subsection title}{size*={28pt}{32pt}}
\setbeamerfont{caption}{size*={16pt}{19pt}}
\setbeamerfont{caption name}{size*={16pt}{19pt}}
\setbeamerfont{itemize/enumerate body}{size*={20pt}{24pt}}
\setbeamerfont{itemize/enumerate subbody}{size*={20pt}{24pt}}
\setbeamerfont{itemize/enumerate subsubbody}{size*={20pt}{24pt}}
\setbeamerfont{itemize item}{size*={20pt}{24pt}}
\setbeamerfont{itemize subitem}{size*={20pt}{24pt}}
\setbeamerfont{itemize subsubitem}{size*={20pt}{24pt}}
\setbeamerfont{description body}{size*={20pt}{24pt}}
\setbeamerfont{description item}{size*={20pt}{24pt}}
\setbeamerfont{quote}{size*={20pt}{24pt}}
\setbeamerfont{quotation}{size*={20pt}{24pt}}
\AtBeginEnvironment{longtable}{\fontsize{20pt}{24pt}\selectfont}
\AtBeginEnvironment{tabular}{\fontsize{20pt}{24pt}\selectfont}
\AtBeginEnvironment{equation}{\fontsize{20pt}{24pt}\selectfont}
\AtBeginEnvironment{equation*}{\fontsize{20pt}{24pt}\selectfont}
\AtBeginEnvironment{align}{\fontsize{20pt}{24pt}\selectfont}
\AtBeginEnvironment{align*}{\fontsize{20pt}{24pt}\selectfont}
\AtBeginEnvironment{itemize}{\fontsize{20pt}{24pt}\selectfont}
\AtBeginEnvironment{enumerate}{\fontsize{20pt}{24pt}\selectfont}
\AtBeginEnvironment{description}{\fontsize{20pt}{24pt}\selectfont}
\AtBeginDocument{\usebeamerfont{normal text}}
\newlength{\AFSlideTableWidth}
% The fixed footer reservation is calibrated with the semantic seven-line regular-body budget.
\setbeamertemplate{footline}{%
  \leavevmode\vbox{%
    \hbox{%
      \begin{beamercolorbox}[wd=\paperwidth,ht=16pt,dp=3pt,center]{author in head/foot}%
        \fontsize{16pt}{19pt}\selectfont Untagged PDF derivative \textbar\ \href{../web/index.html}{HTML reader}%
      \end{beamercolorbox}%
    }%
    \vskip6pt%
  }%
}
% Accessible footer reserved height: 25pt.

\newtheorem{proposition}{Proposition}
\newtheorem{hypothesis}{Hypothesis}
\newtheorem{remark}[theorem]{Remark}
\newtheorem{axiom}{Axiom}
\newtheorem{property}{Property}
\usepackage{bookmark}
\IfFileExists{xurl.sty}{\usepackage{xurl}}{} % add URL line breaks if available
\urlstyle{same}
% fallback for those not using the hyperref driver hyperxmp:
\makeatletter
\@ifundefined{xmpquote}{\newcommand{\xmpquote}[1]{#1}}{}
\makeatother
\hypersetup{
  hidelinks,
  pdfcreator={LaTeX via pandoc}}

\author{\texorpdfstring{}{}}
\date{}

\begin{document}

\begin{frame}{Supplement: hierarchical POMDP methods and parameters}
\protect\phantomsection\label{sec:supp-hierarchical}
This supplement makes the two-level construction of
sec.~14.9 concrete: how location (L1) is
coupled to context (L2) through context-conditioned priors,
\end{frame}

\begin{frame}{Supplement: hierarchical\ldots{} (part 2)}
what the alternating-minimization update actually computes, and the
exact parameters of the executed run.
\end{frame}

\begin{frame}{Supplement: hierarchical\ldots{} (part 3)}
It answers the question the main section brackets --- \emph{by what
mechanism does a second latent level enter the inference at all} --- and
thereby fixes why the hierarchy resolves context above chance while
retaining the separate location-accuracy contrast reported in
sec.~14.9.
\end{frame}

\begin{frame}[fragile]{Generative model for context-gated location
inference}
\protect\phantomsection\label{generative-model-for-context-gated-location-inference}
The two-level POMDP is built by \texttt{build\_hierarchical\_world} in
the \texttt{fedference.pomdp} module.
\end{frame}

\begin{frame}[fragile]{Generative model for\ldots{} (part 2)}
It couples the sentinel's 9-location L1 factor to a 2-state L2 context
factor via \textbf{context-conditioned L1 priors}.

\textbf{L1 (location).} The standard 3x3 grid of
\texttt{build\_sentinel\_world} has \texttt{n\_s\ =} 9 states and sensor
acuity 0.85.
\end{frame}

\begin{frame}[fragile]{Generative model for\ldots{} (part 3)}
\textbf{L2 (context).} A binary \texttt{quiet} / \texttt{alert} state
has a symmetric transition matrix with persistence 0.90 and an initial
uniform prior.
\end{frame}

\begin{frame}[fragile]{Generative model for\ldots{} (part 4)}
\textbf{L1 priors given context.} \texttt{quiet} is uniform over all 9
cells. \texttt{alert} places mass 0.60 at the center cell (the den) and
spreads the residual uniformly.
\end{frame}

\begin{frame}[fragile]{Inference algorithm for top-down empirical
priors}
\protect\phantomsection\label{inference-algorithm-for-top-down-empirical-priors}
The module's \texttt{hierarchical\_infer} function performs 4 passes of
alternating minimization:
\end{frame}

\begin{frame}{Inference algorithm for\ldots{} (part 2)}
\textbf{1. L2 \texorpdfstring{\ensuremath{\to}}{->} L1 empirical prior.} \(\widetilde{\pi}_{0,\mathrm{L1}} =
\sum_c q_{\text{ctx}}[c]\,\pi_{0,\mathrm{L1}\mid c}\), a soft mixture of
the two context-conditioned priors.
\end{frame}

\begin{frame}{Inference algorithm for\ldots{} (part 3)}
\textbf{2. L1 update.} The one-step variational posterior is
\(q_{\text{loc}} = \operatorname{softmax}(\log \widetilde{\pi}_{0,\mathrm{L1}} +
\log A[\text{obs},\,\cdot])\).
\end{frame}

\begin{frame}{Inference algorithm for\ldots{} (part 4)}
\textbf{3. L1 \texorpdfstring{\ensuremath{\to}}{->} L2 marginal evidence.}
\(\ell_c = \log\bigl(\pi_{0,\mathrm{L1}\mid c}^{\top}
A[\text{obs},\,\cdot]\bigr)\), the evidence for context \(c\) from the
observed location likelihood.
\end{frame}

\begin{frame}{Inference algorithm for\ldots{} (part 5)}
\textbf{4. L2 update.}
\(q_{\text{ctx}} = \operatorname{softmax}(\log \pi_{0,\mathrm{L2}} + \ell)\).
\end{frame}

\begin{frame}{Inference algorithm for\ldots{} (part 6)}
After 4 iterations the agent broadcasts both \(q_{\text{loc}}\) and
\(q_{\text{ctx}}\); the colony federates each level independently via a
log-linear pool (eq.~9).
\end{frame}

\begin{frame}{Study parameters for the hierarchical condition}
\protect\phantomsection\label{study-parameters-for-the-hierarchical-condition}
% template-accessible-table-inset
\begingroup
\setlength{\AFSlideTableWidth}{\textwidth}%
\setlength{\LTleft}{\fill}%
\setlength{\LTright}{\fill}%
\begin{longtable}[]{@{}
  >{\raggedright\arraybackslash}p{(\AFSlideTableWidth - 2\tabcolsep) * \real{0.6842}}
  >{\raggedright\arraybackslash}p{(\AFSlideTableWidth - 2\tabcolsep) * \real{0.3158}}@{}}
\label{tbl:hier-params}\tabularnewline
\toprule\noalign{}
\begin{minipage}[b]{\linewidth}\raggedright
Parameter
\end{minipage} & \begin{minipage}[b]{\linewidth}\raggedright
Value
\end{minipage} \\
\midrule\noalign{}
\endfirsthead
\toprule\noalign{}
\begin{minipage}[b]{\linewidth}\raggedright
Parameter
\end{minipage} & \begin{minipage}[b]{\linewidth}\raggedright
Value
\end{minipage} \\
\midrule\noalign{}
\endhead
Agents & 4 \\
Trials & 960 \\
Acuity & 0.85 \\
\bottomrule\noalign{}
\end{longtable}
\endgroup
\end{frame}

\begin{frame}{Study parameters for the\ldots{} (part 2)}
The executed hierarchical configuration is summarized in
tbl.~17.
\end{frame}

\end{document}
